% Chapter 5. The ErisML compiler.
% Author: agent C. Every claim is sourced from the erisml-compiler worktree at
% C:\source\erisml-compiler-monitor (version 0.9.0, src/erisml_compiler/__init__.py:11).
% Paths in "% src:" comments are relative to that worktree unless stated otherwise.

\chapter{The ErisML compiler}\label{ch:compiler}

The ErisML compiler turns a description of a moral situation into a typed graph, reads that graph
through several ethical theories, and records every step in a hash chain. Its input is either
structured ErisML source or natural-language text. Its output is an intermediate representation
that downstream consumers can audit, replay and contest. This chapter describes the compiler as
the code implements it in release 0.9.0. Where the design documents and the code differ, the
chapter reports what the code does and names the difference.

The chapter follows the data. Section~\ref{sec:compiler-split} gives the two-layer split between a
descriptive substrate and the theory-bound projections. Section~\ref{sec:compiler-pipeline} walks
the pass pipeline. Sections~\ref{sec:compiler-extraction} and~\ref{sec:compiler-canon} cover
extraction, the three tiers, canonicalization and human correction.
Section~\ref{sec:compiler-projections} covers the four projections and the disagreement record.
Section~\ref{sec:compiler-emdag} treats the ethical-module graph in depth.
Section~\ref{sec:compiler-fsm} covers the finite-state layer and the hardware target.
Section~\ref{sec:compiler-audit} covers the audit chain and the training export.
Section~\ref{sec:compiler-boundaries} closes with the trust boundaries and the limits the
architecture states about itself.

% src: src/erisml_compiler/__init__.py:11 (__version__ = "0.9.0")
% src: docs/architecture.md:1-11

\section{A descriptive substrate sits below theory-bound projections}\label{sec:compiler-split}

The compiler separates two layers. The lower layer is the moral substrate. It is a typed directed
graph, the \code{MoralGraph}, that records who did what to whom, under what authority and under
what standing commitments. The upper layer is a set of projections. Each projection reads the
same substrate and emits the kind of object its ethical theory produces.
% src: docs/architecture.md:46-73

The substrate has six node kinds and eleven edge kinds. The node kinds are stakeholder, act,
maxim, commitment, fact and norm. A maxim node holds an action under a particular description,
which is the unit a Kantian test operates on. The edge kinds are listed in
Table~\ref{tab:compiler-edges}. Node and edge payloads carry the original schema objects verbatim,
so the graph view loses nothing that the flat lists held.
% src: docs/architecture.md:75-104
% src: src/erisml_compiler/projections/substrate.py:42-67 (Maxim)

\begin{table}[t]
\centering\footnotesize
\begin{tabular}{lll}
\hline
Edge kind & Direction & Payload \\
\hline
\code{performs} & stakeholder to act & \\
\code{imposes\_on} & act to stakeholder & severity, confidence \\
\code{consents\_to} & stakeholder to act & given, under\_duress, informed \\
\code{holds\_commitment} & stakeholder to commitment & \\
\code{commitment\_binds} & commitment to beneficiary & \\
\code{treats\_as} & act to stakeholder & role in \{end, means, mere\_means\} \\
\code{under\_maxim} & act to maxim & \\
\code{coerces} & stakeholder to stakeholder & \\
\code{surfaces\_fact} & act to fact & \\
\code{fact\_subject} & fact to stakeholder & \\
\code{would\_violate\_if\_universalised} & act to norm & \\
\hline
\end{tabular}
\caption{The eleven typed edges of the moral graph. Each edge has a fixed source and target node
kind, and three edges carry a payload that later passes read.}
\label{tab:compiler-edges}
\end{table}
% src: docs/architecture.md:88-100

Four projections ship. The consequentialist projection produces a per-stakeholder tensor, a
verdict from the DEME evaluator and distributive statistics. The deontic projection produces four
categorical gate findings. The virtue projection and the care projection each produce three
findings. Every projection result carries the same minimum fields, so the orchestrator can compare
them without knowing their internals.
% src: docs/architecture.md:156-163
% src: src/erisml_compiler/projections/base.py:94-125

The split exists so that the compiler never aggregates across theories in silence. The
architecture document states the reason directly. Aggregation across theories is itself a
metaethical choice, and the compiler leaves that choice to its caller.
Figure~\ref{fig:compiler-two-layer} shows the two layers and the comparison that joins them.
% src: docs/architecture.md:68-73

\begin{figure}[t]
\centering
% Chapter 5. The two-layer split: one descriptive substrate, four theory-bound projections.
% src: C:\source\erisml-compiler-monitor\docs\architecture.md:46-73, 156-186
% src: C:\source\erisml-compiler-monitor\src\erisml_compiler\pipeline\orchestrator.py:437-454
\begin{tikzpicture}[house]
% the substrate
\node[panel=Teal, minimum width=128mm, minimum height=27mm] (sub) {};
\node[stepname=Teal, anchor=north] at ($(sub.north)+(0,-2mm)$) {Moral substrate, a typed MoralGraph};
\node[desc, anchor=north] at ($(sub.north)+(0,-6.5mm)$) {descriptive, who did what to whom, under what authority and commitments};
\node[detail, text width=36mm] at ($(sub.south)+(-42mm,6.5mm)$) {6 node kinds\\stakeholder, act, maxim,\\commitment, fact, norm};
\node[detail, text width=38mm] at ($(sub.south)+(0,6.5mm)$) {11 typed edge kinds\\imposes\_on, consents\_to,\\treats\_as, coerces, \dots};
\node[detail, text width=36mm] at ($(sub.south)+(42mm,6.5mm)$) {canonical SHA-256\\$h_G$ into the\\audit record};
% the four projections
\node[panel=Green, text width=25mm, minimum height=24mm, below=11mm of sub.south west, anchor=north west] (con)
  {\textcolor{hsGreen}{\scriptsize\textbf{Consequentialist}}\\[1pt]{\scriptsize EM-DAG, tensor,\\DEME verdict, Gini}\\[1pt]{\scriptsize permit, escalate}};
\node[panel=Blue, text width=25mm, minimum height=24mm, right=2mm of con] (deo)
  {\textcolor{hsBlue}{\scriptsize\textbf{Deontic}}\\[1pt]{\scriptsize 4 gates, Z3\\universalizability}\\[1pt]{\scriptsize permit, forbid, escalate}};
\node[panel=Purple, text width=25mm, minimum height=24mm, right=2mm of deo] (vir)
  {\textcolor{hsPurple}{\scriptsize\textbf{Virtue}}\\[1pt]{\scriptsize 3 findings, character,\\commitments, asymmetry}\\[1pt]{\scriptsize permit, escalate}};
\node[panel=Orange, text width=25mm, minimum height=24mm, right=2mm of vir] (car)
  {\textcolor{hsOrange}{\scriptsize\textbf{Care}}\\[1pt]{\scriptsize 3 findings, relations,\\dependence}\\[1pt]{\scriptsize permit, escalate}};
\foreach \p in {con,deo,vir,car} { \draw[flow] (sub.south -| \p.north) -- (\p.north); }
% comparison and record
\node[panel=Amber, text width=44mm, anchor=north] (cmp) at ($(deo.south)!0.5!(vir.south)+(0,-10mm)$)
  {\textcolor{hsAmber}{\textbf{Compare polarities}}\\[1pt]{\scriptsize more than one non-neutral polarity}};
\node[panel=Grey, text width=34mm, right=5mm of cmp] (rec)
  {\textcolor{hsGrey}{\textbf{Disagreement record}}\\[1pt]{\scriptsize all verdicts, no winner,\\choice left to the caller}};
\foreach \p in {con,deo,vir,car} { \draw[flow] (\p.south) -- (cmp.north); }
\draw[flow] (cmp) -- (rec);
\end{tikzpicture}

\caption{One descriptive substrate read by four theory-bound projections. Each projection lists
the polarities its shipped code can emit, and the comparison writes a record of disagreement
instead of a combined verdict.}
\label{fig:compiler-two-layer}
\end{figure}
% src: src/erisml_compiler/projections/{consequentialist,deontic,virtue,care_ethics}.py (verdict rules)

\section{The pipeline runs twelve numbered passes with one inserted graph pass}\label{sec:compiler-pipeline}

The orchestrator function \code{compile\_document} runs the pipeline. The command-line entry point
\code{cmd\_compile} calls it. Each pass runs inside a context manager that times it and appends a
\code{PassRecord} with its index, its name, the tier and its duration in milliseconds. The list of
pass records becomes part of the audit record.
% src: src/erisml_compiler/pipeline/orchestrator.py:269-509
% src: src/erisml_compiler/audit/provenance.py:11-25
% src: src/erisml_compiler/ir/schemas.py:355-362

Figure~\ref{fig:compiler-passes} shows the passes as the orchestrator records them. Pass~0 loads
the input. A structured input goes straight to the graph pass. A text input goes through
segmentation, extraction and canonicalization first. Pass~7.5 builds the typed graph and its hash.
Pass~8 runs every enabled projection. Pass~10 is kept as an empty record for chain continuity,
because the DEME verdict is now computed inside the consequentialist projection. Pass~11 detects
conflicts between dimensions of the moral vector. Pass~12 finalizes the audit record. Two optional
steps then run after the audit, a strategic analysis and a decision proof.
% src: src/erisml_compiler/pipeline/orchestrator.py:281-507

\begin{figure}[t]
\centering
% Chapter 5. The compiler passes as the orchestrator records them (pass index in each circle).
% src: C:\source\erisml-compiler-monitor\src\erisml_compiler\pipeline\orchestrator.py:269-509
\begin{tikzpicture}[house,
  pp/.style={panel=#1, text width=32mm, minimum height=19mm}]
% row 1, left to right
\node[pp=Blue] (p0) at (0,0) {\textcolor{hsBlue}{\textbf{Ingestion}}\\[1pt]{\scriptsize text loader, or the Tier 1\\structured loader}};
\node[pp=Blue] (p1) at (42mm,0) {\textcolor{hsBlue}{\textbf{Segmentation}}\\[1pt]{\scriptsize paragraphs tagged as\\dialogue, norm, decision}};
\node[pp=Purple] (p2) at (84mm,0) {\textcolor{hsPurple}{\textbf{Extraction}}\\[1pt]{\scriptsize mock, rule, probe or LLM\\with an optional critic}};
% row 2, right to left
\node[pp=Purple] (p7) at (84mm,-29mm) {\textcolor{hsPurple}{\textbf{Canonicalization}}\\[1pt]{\scriptsize Jaccard registry or\\LaBSE cosine}};
\node[pp=Teal] (p75) at (42mm,-29mm) {\textcolor{hsTeal}{\textbf{Graph identity}}\\[1pt]{\scriptsize promote to MoralGraph,\\compute $h_G$}};
\node[pp=Green] (p8) at (0,-29mm) {\textcolor{hsGreen}{\textbf{Projections}}\\[1pt]{\scriptsize four theories, EM-DAG\\inside the first}};
% row 3, left to right
\node[pp=Grey] (p10) at (0,-58mm) {\textcolor{hsGrey}{\textbf{DEME record}}\\[1pt]{\scriptsize empty, kept for chain\\continuity}};
\node[pp=Amber] (p11) at (42mm,-58mm) {\textcolor{hsAmber}{\textbf{Conflict detection}}\\[1pt]{\scriptsize opposite signs with\\magnitude above 0.5}};
\node[pp=Orange] (p12) at (84mm,-58mm) {\textcolor{hsOrange}{\textbf{Audit record}}\\[1pt]{\scriptsize $h_{\mathrm{IR}}$, $h_G$, profiles,\\pass timings}};
% step circles carry the pass index the code records
\foreach \n/\c/\t in {p0/Blue/0, p1/Blue/1, p2/Purple/2, p7/Purple/7, p75/Teal/7.5, p8/Green/8, p10/Grey/10, p11/Amber/11, p12/Orange/12}
  { \node[step=\c] at (\n.north west) {\t}; }
% flow
\draw[flow] (p0) -- node[flowlabel, above]{text} (p1);
\draw[flow] (p1) -- (p2);
\draw[flow] (p2) -- (p7);
\draw[flow] (p7) -- (p75);
\draw[flow] (p75) -- (p8);
\draw[flow] (p8) -- (p10);
\draw[flow] (p10) -- (p11);
\draw[flow] (p11) -- (p12);
\draw[flow] (p0.south) -- ++(0,-4.5mm) -| node[flowlabel, pos=0.3, above]{Tier 1} ([xshift=-6mm]p75.north);
\draw[untrusted] ($(p2.north)+(0,9mm)$) node[flowlabel, above]{Tier 3 model output} -- (p2.north);
% after the audit
\node[panel=Grey, text width=64mm, below=6mm of p12.south east, anchor=north east] (post)
  {{\scriptsize after Pass 12, strategic analysis and the decision proof}};
\draw[flow] (p12) -- (p12 |- post.north);
\end{tikzpicture}

\caption{The compiler passes in the order the orchestrator records them, with the recorded pass
index in each circle. A structured input skips segmentation, extraction and canonicalization, and
only Tier 3 brings model-written content into the pipeline.}
\label{fig:compiler-passes}
\end{figure}
% src: src/erisml_compiler/pipeline/orchestrator.py:282-356 (passes 0, 1, 2, 7, 75)
% src: src/erisml_compiler/pipeline/orchestrator.py:363-480 (passes 8, 10, 11, 12)
% src: src/erisml_compiler/pipeline/orchestrator.py:485-507 (post-audit steps)

The numbering carries history. The orchestrator docstring lists twelve passes from an earlier
specification, with tensorization at Pass~8 and DEME evaluation at Pass~10. The architecture
document lists the EM-DAG at Pass~8 and the projections at Pass~8.5. The code records the
projection pass under index~8 and records Pass~7.5 under the integer index~75. The EM-DAG runs
inside the projection pass. This chapter uses the code's numbering.
% src: src/erisml_compiler/pipeline/orchestrator.py:1-17 (docstring pass map)
% src: docs/architecture.md:188-206 (pass table)
% src: src/erisml_compiler/pipeline/orchestrator.py:350 (record_pass(passes, 75, "graph_identity", ...))
% src: src/erisml_compiler/pipeline/orchestrator.py:363 (record_pass(passes, 8, "projection_pass", ...))

The options that shape a run are collected in \code{CompileOptions}. They name the tier, the
extractor, an optional critic extractor, the EM-DAG profile, the set of projections, an optional
ethos profile, the canonicalizer and the tensor rank between one and six. One option,
\code{strict\_v3}, governs a fallback. By default, if the DEME V3 tensor path in \code{erisml-lib}
is missing or raises, the compiler falls back to an older tensor builder and logs a warning. With
\code{strict\_v3} set, the same failure raises \code{StrictV3Error} instead. The option exists so
that a research run cannot downgrade without notice.
% src: src/erisml_compiler/pipeline/orchestrator.py:122-167 (CompileOptions)
% src: src/erisml_compiler/pipeline/orchestrator.py:53-119 (StrictV3Error, _produce_v3_tensor)

\section{Three tiers feed one evaluator}\label{sec:compiler-extraction}

\subsection{The tiers differ only in how structure enters}

The compiler defines three tiers in \code{tiers.py}. Tier~1, named \code{GEOMETRIC}, reads
pre-parsed structured input and runs the finite-state machines and the evaluator. The code
describes it as bounded in memory, deterministic in dispatch and feasible in fixed point. Tier~2,
named \code{RULES}, extracts structure from natural-language text with patterns and keywords. It
runs offline and deterministically. Tier~3, named \code{LLM}, extracts structure with a language
model. It is the most expressive tier and runs in batch only. All three tiers emit the same
intermediate representation, and the evaluator core is shared.
% src: src/erisml_compiler/tiers.py:1-23
% src: src/erisml_compiler/tiers.py:31-37

The tier also decides two properties. Only Tier~1 is marked as castable to silicon, because only
its full pipeline avoids text processing. Only Tier~3 requires network access. The tier is
detected from the file extension when the caller does not name it. Files ending in
\code{.json}, \code{.erisml}, \code{.yaml} or \code{.yml} select Tier~1. Files ending in
\code{.txt} or \code{.md} select Tier~2. Tier~3 is never selected automatically.
% src: src/erisml_compiler/tiers.py:39-73

The Tier~1 loader reads JSON or ErisML source. ErisML source is the YAML form that the code
generator emits, with the representation's own keys for stakeholders, relations, events,
commitments, norms, ethical facts and an open \code{extra} block. The loader records the SHA-256
of the raw file in the document record. The home-care twin of Part~III uses this loader to read its
scene files.
% src: src/erisml_compiler/ingestion/structured_loader.py:1-23, 48-76
% src: C:\source\gtc-prototype\twin\brain.py:301-319 (load_structured_input on the scene)

\subsection{Extractors share one result type}

Every extractor implements one abstract method, \code{extract}, which takes a document and its
segments and returns an \code{ExtractorResult}. The result holds lists of stakeholders, relations,
events, commitments, norms, ethical facts and conflicts, a canonical-form tag, free metadata and an
optional graph. When an extractor fills the graph field, the orchestrator skips the promotion step
at Pass~7.5 and treats the extractor's graph as the source of truth.
% src: src/erisml_compiler/annotation/base.py:26-68
% src: src/erisml_compiler/pipeline/orchestrator.py:345-356

Four extractors ship. The mock extractor returns hand-written fixtures for three example texts,
selected by the SHA-256 of the text, and raises an error on any other document. The rule extractor
is the production Tier~2 extractor. It uses regular expressions for commitment verbs, coercion,
deception, collective targets and profession nouns. The probe extractor is labelled Tier~2.5. It
encodes each segment with a frozen LaBSE sentence encoder~\cite{feng2022labse} and classifies it
with trained probe heads into a stakeholder role and an ethical-fact kind. Without a trained
checkpoint it returns a fixed low-confidence default, and its metadata records that no checkpoint
was loaded. The LLM extractor issues four prompts per document, for stakeholders, commitments,
ethical facts and the canonical form. It takes the first balanced JSON array from each reply,
validates each element against the schema and keeps the elements that pass.
% src: src/erisml_compiler/annotation/mock_extractor.py:1-10
% src: src/erisml_compiler/annotation/rule_extractor.py:1-60
% src: src/erisml_compiler/annotation/probe_extractor.py:1-20, 99-185
% src: src/erisml_compiler/annotation/llm_extractor.py:1-25

Segmentation precedes extraction for text input. The segmenter splits on blank lines and tags
each paragraph as dialogue, norm statement, decision point or background by keyword patterns.
% src: src/erisml_compiler/segmentation/segmenter.py:1-37

\subsection{A critic extractor measures agreement and marks what it cannot corroborate}

A second extractor can run as a critic over the first. The critic pass computes four agreement
scores between the primary result $P$ and the critic result $C$. Three are Jaccard overlaps,
\begin{equation}
J(A,B) = \frac{|A \cap B|}{|A \cup B|}, \qquad J(\emptyset,\emptyset) = 1,
\label{eq:compiler-jaccard}
\end{equation}
taken over stakeholder identifiers, commitment identifiers and ethical-fact kinds. The fourth is
an indicator $a_{\mathrm{cf}} \in \{0,1\}$ that both extractors chose the same non-empty
canonical form. The overall agreement is the mean of the four,
\begin{equation}
\bar{a} = \tfrac{1}{4}\bigl(J_{\mathrm{stake}} + J_{\mathrm{commit}} + J_{\mathrm{kind}} + a_{\mathrm{cf}}\bigr).
\label{eq:compiler-agreement}
\end{equation}
Each stakeholder that the critic did not find is marked \code{requires\_review}. When $\bar{a}$
falls below a review threshold, 0.5 by default, a document-level review note is added. The primary
result is kept as the output. If the critic raises an error, the primary result passes through
with the error recorded in its metadata.
% src: src/erisml_compiler/annotation/critic.py:21-50 (CriticReport, _set_overlap)
% src: src/erisml_compiler/annotation/critic.py:56-142 (critic_pass)
% src: src/erisml_compiler/annotation/critic.py:145-179 (CriticExtractor)

\section{Canonicalization and correction}\label{sec:compiler-canon}

\subsection{Two canonicalizers map a situation to a registered form}

Canonicalization assigns the situation a stable tag from a registry of canonical forms. The
registry lives in \code{ontology/canonical\_forms.yaml}. The orchestrator builds a one-paragraph
summary from the document title, the first 400 characters of text, up to eight ethical-fact
descriptions and up to four commitments. It passes the summary and the registry to a
canonicalizer.
% src: src/erisml_compiler/pipeline/orchestrator.py:242-266
% src: src/erisml_compiler/pipeline/orchestrator.py:327-343

Two canonicalizers share one interface. The registry canonicalizer tokenizes the summary and each
form's description and tag into lowercase words longer than two characters, removes stopwords, and
scores each form by the Jaccard similarity of Equation~\eqref{eq:compiler-jaccard}. It returns the
best form when the score reaches 0.10. The embedding canonicalizer encodes the summary and each
form's text with LaBSE, normalizes the embeddings and scores by the dot product, which equals the
cosine similarity $\cos(\mathbf{e}_s, \mathbf{e}_f)$ for unit vectors. It returns the best form
when the score reaches 0.55. Both thresholds can be changed, the second through an environment
variable. When the best score is below threshold, the result carries no tag, and the extractor's
own tag is kept.
% src: src/erisml_compiler/canonicalizer/registry.py:68-120 (threshold 0.10)
% src: src/erisml_compiler/canonicalizer/labse.py:20-110 (threshold 0.55, ERISML_LABSE_THRESHOLD)
% src: src/erisml_compiler/pipeline/orchestrator.py:335-336

The selector \code{auto\_canonicalizer} prefers the embedding canonicalizer and falls back to the
registry canonicalizer when the sentence-encoder package cannot load. Every result records its
backend, its score, whether the tag is a registered form and a list of evidence strings. The
orchestrator stores all of these under \code{ir.extra["canonicalization"]}.
% src: src/erisml_compiler/canonicalizer/base.py:8-71
% src: src/erisml_compiler/pipeline/orchestrator.py:337-343

The home-care twin uses the registry canonicalizer, not the embedding one. Its classifier
canonicalizes, and its chooser sees only the canonical state.
% src: C:\source\gtc-prototype\twin\brain.py:315-319

\subsection{Human corrections are patches with hashes on both sides}

The correction package closes a loop from human review back into training data. A corrections
file names a corrector, a rationale and a list of patches. Each patch is one of three operations,
\code{set}, \code{add} or \code{remove}, applied at a dotted path such as
\code{stakeholders.village.type}. Paths can address a scalar field, a whole entity in a keyed
collection, or a field inside an entity. The keyed collections are stakeholders, commitments,
events, ethical facts, conflicts, norms and relations.
% src: src/erisml_compiler/correction/__init__.py:1-13
% src: src/erisml_compiler/correction/corrector.py:1-46, 72-173

Applying corrections proceeds in a fixed order. The corrector hashes the representation, applies
each patch, and counts the patches that fail without stopping. It then rebuilds the
representation through schema validation and hashes it again. A structural diff, organised per
entity, summarizes the change. The \code{CorrectionRecord} holds the corrector, the rationale, the
time, the counts of applied and failed patches, the hashes before and after, and the diff summary.
The record is appended to \code{ir.extra["corrections"]}, so the history of corrections travels
with the representation.
% src: src/erisml_compiler/correction/corrector.py:176-235
% src: src/erisml_compiler/correction/diff.py:1-40

\section{Four projections, and disagreement as a result}\label{sec:compiler-projections}

\subsection{The substrate view}

Each projection reads a \code{MoralSubstrate}. When the graph is present, the substrate is
derived from the graph. Otherwise it is derived from the flat lists. Besides the lists, the
substrate holds three derived objects. The first is the maxim, with its agent, a coarse action
kind, a polarity of affirmed or negated, and a purpose. The second is the list of consent states.
The third is the list of authority legitimacies.
% src: src/erisml_compiler/projections/substrate.py:135-164
% src: src/erisml_compiler/projections/substrate.py:42-67

\subsection{The consequentialist projection}

The consequentialist projection runs the EM-DAG of Section~\ref{sec:compiler-emdag}, projects
the module outputs onto the moral vector, builds the moral tensor and asks the DEME bridge for a
verdict. It returns the per-party verdicts and the Gini and worst-off statistics that the DEME
verdict carries. Its metadata records the DAG name, whether an ethos profile was applied and a
count of graph nodes and edges by kind. The orchestrator copies this projection's outputs into
the older top-level fields, so consumers written before the split keep working.
% src: src/erisml_compiler/projections/consequentialist.py:232-357
% src: src/erisml_compiler/pipeline/orchestrator.py:373-410

The DEME bridge is a cascade of fixed comparisons over the vector. It returns
\code{tragic\_conflict\_escalate} when the externality score is at most $-0.85$ and some
commitment is active but defeasible. It returns \code{prohibited} when harm is below $-0.85$,
care is below $0.5$ and legitimacy is below zero. Further branches return
\code{permitted\_with\_residue} and \code{permitted}. Chapter~6 treats DEME in full.
% src: src/erisml_compiler/erisml_backend/deme_bridge.py:40-120

\subsection{The deontic projection runs four gates and a solver}

The deontic projection emits four \code{GateFinding} records. A gate finding is categorical. It
has a name, a flag \code{passed}, a reason tied to substrate elements, a severity from minor to
catastrophic and a list of implicated stakeholders.
% src: src/erisml_compiler/projections/base.py:32-57
% src: src/erisml_compiler/projections/deontic.py:45-91

The \emph{universalizability} gate tests the maxim's action kind. When the Z3 solver~\cite{demoura2008z3}
is installed and the action kind is in its model, the gate runs a satisfiability check. The model
has nine Boolean institutional facts, among them \code{truth\_telling\_default},
\code{promises\_create\_trust} and \code{non\_consenting\_party\_standing}. For each action kind
the model lists the facts that universalizing the maxim destroys, $D$, and the facts the act
presupposes, $P$. The contradiction-in-conception test asks whether
\begin{equation}
\Phi_{\mathrm{CIC}} = \bigwedge_{d \in D} \neg d \;\wedge\; \bigwedge_{p \in P} p
\label{eq:compiler-cic}
\end{equation}
is satisfiable. If it is not, the maxim fails with a contradiction in conception. The model of
deception, for example, destroys and presupposes the same fact, \code{truth\_telling\_default}, so
$\Phi_{\mathrm{CIC}}$ is unsatisfiable. If $\Phi_{\mathrm{CIC}}$ is satisfiable, the solver adds
the agent's own ends $E$, four facts that a rational agent needs, and checks
\begin{equation}
\Phi_{\mathrm{CIW}} = \Phi_{\mathrm{CIC}} \;\wedge\; \bigwedge_{e \in E} e .
\label{eq:compiler-ciw}
\end{equation}
If $\Phi_{\mathrm{CIW}}$ is unsatisfiable, the maxim fails with a contradiction in will. Coercion
and refusal of help fail this way. Otherwise the gate passes, and the satisfying assignment is
recorded in the finding for audit. A negated maxim, an action kind outside the model, or a missing
solver all route to a lookup table of hand-written dependencies instead. The finding records which
path ran.
% src: src/erisml_compiler/delta/universalizability_smt.py:1-58 (fact universe, agent ends)
% src: src/erisml_compiler/delta/universalizability_smt.py:107-196 (_ZModel, _Z3_ACTION_MODEL)
% src: src/erisml_compiler/delta/universalizability_smt.py:218-367 (test_universalizability_smt, _solve)
% src: src/erisml_compiler/projections/deontic.py:95-168

The solver adds consistency checking, not moral content. The source file says so. The fact
universe and each action's destroyed and presupposed facts are written by hand. What the solver
contributes is that the hand-written model cannot contradict itself without the contradiction
showing.
% src: src/erisml_compiler/delta/universalizability_smt.py:46-58

The \emph{mere means} gate fires on any \code{treats\_as} edge whose role is \code{mere\_means}.
Without such edges it falls back to the maxim's own record, and then to any stakeholder whose
consent was not given. The \emph{valid consent} gate fires when any consent state is absent,
under duress or uninformed. The \emph{legitimate authority} gate fires when any authority in the
substrate is not legitimate. The first three gates carry severity grave and the fourth moderate.
The verdict follows from the failures,
\begin{equation}
v_{\mathrm{deon}} =
\begin{cases}
\code{forbidden} & \text{if a grave or catastrophic gate fails,}\\
\code{requires\_review} & \text{else if a moderate gate fails,}\\
\code{permissible} & \text{otherwise.}
\end{cases}
\label{eq:compiler-deontic}
\end{equation}
% src: src/erisml_compiler/projections/deontic.py:71-79 (verdict rule)
% src: src/erisml_compiler/projections/deontic.py:172-296 (other three gates)

\subsection{The virtue and care projections}

The virtue projection emits three findings. The character finding maps the action kind to a
virtue and vice pair, such as \code{honesty} and \code{deception} or \code{fidelity} and \code{perfidy}, from a table of six
action kinds. Deceiving and imposing an externality express the vice. Negation flips the reading,
so refraining from a vice expresses the virtue. The commitment finding reports the count of
standing commitments as context and always passes. The asymmetry finding fires when the act
imposes on a stakeholder labelled vulnerable, patient or dependent. The verdict is \code{vicious}
when two or more grave or catastrophic findings fail, \code{requires\_practical\_wisdom} when any
finding fails, and \code{virtuous} otherwise. In the shipped code all three findings carry
moderate or minor severity, so the projection's verdict is \code{virtuous} or
\code{requires\_practical\_wisdom}.
% src: src/erisml_compiler/projections/virtue.py:36-46 (_ACTION_VIRTUE_AXES)
% src: src/erisml_compiler/projections/virtue.py:53-90 (verdict rule)
% src: src/erisml_compiler/projections/virtue.py:92-228 (findings and severities)

The care projection follows the care-ethics tradition, which reasons from relations and dependence
rather than from isolated acts. It emits three findings. The relational-attentiveness finding
fails when the substrate holds no relations. The asymmetric-responsibility finding fails when a
dependent, a patient or a highly vulnerable party is present. It flags an asymmetry and does not
judge whether the duty was met. The dependency-response finding fails, with grave severity, when an
\code{imposes\_on} edge targets a dependent party. The verdict is \code{uncaring} when two or more
grave findings fail, \code{requires\_caring\_attention} when any finding fails and \code{caring}
otherwise. Only the dependency-response finding is grave, so in the shipped code the
\code{uncaring} verdict cannot occur. Both the virtue and the care projections can therefore
permit or escalate, and neither can forbid.
% src: src/erisml_compiler/projections/care_ethics.py:1-27
% src: src/erisml_compiler/projections/care_ethics.py:39-213

\subsection{Disagreement is recorded, never resolved}

Each projection's native verdict maps to one of four polarities, \code{permit}, \code{forbid},
\code{escalate} and \code{neutral}. Comparing polarities keeps vocabulary from looking like
disagreement, since \code{permitted}, \code{permissible}, \code{virtuous} and \code{caring} all map
to \code{permit}. Let $\pi_k$ be the polarity of projection $k$. The orchestrator writes a
disagreement record when
\begin{equation}
\bigl|\{\pi_k\}_k \setminus \{\code{neutral}\}\bigr| > 1 .
\label{eq:compiler-disagree}
\end{equation}
The record holds every native verdict, every polarity and a fixed note. The note says that the
compiler surfaces all verdicts without aggregating and leaves the choice to the caller. The
compiler does not name a winner.
% src: src/erisml_compiler/projections/base.py:60-91 (VerdictPolarity, _DEFAULT_POLARITY_MAP)
% src: src/erisml_compiler/pipeline/orchestrator.py:437-454
% src: docs/architecture.md:165-186

The polarity map has a consequence that a reader of the code should know. The map lists
\code{permitted}, \code{requires\_human\_review}, \code{tragic\_conflict\_escalate} and
\code{forbidden\_action} for the consequentialist projection. The DEME bridge also returns
\code{prohibited}, \code{permitted\_with\_residue} and \code{indeterminate}. Those three strings are
not in the map, so they map to \code{neutral}. A consequentialist \code{prohibited} verdict
therefore drops out of the comparison in Equation~\eqref{eq:compiler-disagree}.
% src: src/erisml_compiler/projections/base.py:67-91 (map entries)
% src: src/erisml_compiler/erisml_backend/deme_bridge.py:61-120 (verdict strings returned)
% src: src/erisml_compiler/projections/consequentialist.py:336-340 (polarity not set explicitly)

\section{The EM-DAG}\label{sec:compiler-emdag}

The EM-DAG is the graph of ethical modules that feeds the consequentialist projection. Each module
evaluates one ethical concern and writes one dimension of the moral vector. Edges run from a module
to the modules that read its output. The topology is the DEME profile, so a different deployment
can load a different graph.
% src: src/erisml_compiler/em_dag/__init__.py:1-13
% src: docs/architecture.md:228-254

\subsection{The module contract}

A module subclasses \code{EthicalModule}. It declares three class attributes, a \code{name}, the
moral-vector \code{dimension} it writes and a tuple of \code{dependencies}. The class refuses to be
defined without a name and a dimension. It implements one method,
\code{evaluate(ir, upstream)}, which receives the full representation and the outputs of its
declared dependencies, and returns an \code{EMOutput}.
% src: src/erisml_compiler/em_dag/base.py:10-51

An \code{EMOutput} holds the module name, a \code{DimensionScore}, the identifiers of the facts
that contributed and the upstream dependencies. A dimension score has a value $v \in [-1,1]$, a
confidence $c \in [0,1]$, an uncertainty $u \in [0,1]$, a direction of positive, negative or
neutral, a list of source spans and an explanation. Positive values mean the dimension is
respected. Negative values mean it is violated.
% src: src/erisml_compiler/ir/schemas.py:264-272 (DimensionScore)
% src: src/erisml_compiler/em_dag/modules/_helpers.py:1-6 (sign convention)

\subsection{Validation and evaluation}

The \code{EMDAG} class validates on construction. Every declared dependency must name a module in
the graph, or construction raises. It then orders the modules with Kahn's
algorithm~\cite{kahn1962topological}. Each module starts with an in-degree equal to its number of
dependencies. Modules with in-degree zero enter a queue. Removing a module from the queue
decrements the in-degree of its dependants and enqueues those that reach zero. If the resulting
order is shorter than the module set, the graph has a cycle, and construction raises with the
modules still holding positive in-degree.
% src: src/erisml_compiler/em_dag/dag.py:67-113

Evaluation walks the stored order. For each module it collects the outputs of that module's
dependencies, which are already computed, and calls \code{evaluate}. A module therefore never
sees the output of a module it did not declare. The result is a map from module name to output.
% src: src/erisml_compiler/em_dag/dag.py:133-144

\subsection{The default profile}

A profile is a YAML file with a name and a list of modules, each given as a
\code{module.path:ClassName} reference. The loader imports each class and instantiates it with no
arguments. The default profile, \code{default\_em\_dag\_v0.1}, lists ten modules in three groups.
Five are roots with no dependencies. Four form a second tier. One composes the others.
Figure~\ref{fig:compiler-emdag} shows the graph as configured.
% src: src/erisml_compiler/em_dag/dag.py:152-179 (load_profile)
% src: src/erisml_compiler/em_dag/profiles/default.yaml:1-24

\begin{figure}[t]
\centering
% Chapter 5. The default EM-DAG as configured (default.yaml), circles give the topological order.
% src: C:\source\erisml-compiler-monitor\src\erisml_compiler\em_dag\profiles\default.yaml:11-24
% src: C:\source\erisml-compiler-monitor\src\erisml_compiler\em_dag\modules\*.py (dependencies)
% src: C:\source\erisml-compiler-monitor\src\erisml_compiler\silicon\examples\erisml_em_dag.cpp:5 (order)
\begin{tikzpicture}[house,
  em/.style={panel=#1, text width=29mm, minimum height=13mm}]
% roots
\node[em=Teal] (harm) at (-6mm,4.5mm) {\textcolor{hsTeal}{\textbf{harm}}\\{\scriptsize\ttfamily physical\_harm}};
\node[em=Teal] (legit) at (-6mm,-33.5mm) {\textcolor{hsTeal}{\textbf{legitimacy}}\\{\scriptsize\ttfamily legitimacy\_trust}};
\node[em=Teal] (rights) at (-6mm,-64mm) {\textcolor{hsTeal}{\textbf{rights}}\\{\scriptsize\ttfamily rights\_respect}};
\node[em=Teal] (fair) at (44mm,-64mm) {\textcolor{hsTeal}{\textbf{fairness}}\\{\scriptsize\ttfamily fairness\_equity}};
\node[em=Teal] (epi) at (88mm,-64mm) {\textcolor{hsTeal}{\textbf{epistemic}}\\{\scriptsize\ttfamily epistemic\_quality}};
% second tier
\node[em=Blue] (ext) at (44mm,14mm) {\textcolor{hsBlue}{\textbf{externality}}\\{\scriptsize\ttfamily third\_party\_externality}};
\node[em=Blue] (care) at (44mm,-5mm) {\textcolor{hsBlue}{\textbf{care}}\\{\scriptsize\ttfamily care\_protection}};
\node[em=Blue] (auto) at (44mm,-24mm) {\textcolor{hsBlue}{\textbf{autonomy}}\\{\scriptsize\ttfamily autonomy\_consent}};
\node[em=Blue] (fid) at (44mm,-43mm) {\textcolor{hsBlue}{\textbf{fidelity}}\\{\scriptsize\ttfamily vow\_fidelity}};
% composition
\node[em=Orange] (rep) at (88mm,-14.5mm) {\textcolor{hsOrange}{\textbf{repair}}\\{\scriptsize\ttfamily repair\_residue}};
% topological order from Kahn's algorithm
\foreach \n/\c/\k in {harm/Teal/1, rights/Teal/2, fair/Teal/3, legit/Teal/4, epi/Teal/5, ext/Blue/6, care/Blue/7, auto/Blue/8, fid/Blue/9, rep/Orange/10}
  { \node[step=\c] at (\n.north west) {\k}; }
% dependency edges that carry data
\draw[flow] (harm.east) -- (care.west);
\draw[flow] (legit.east) -- (auto.west);
\draw[flow] (ext.east) -| ([xshift=-4mm]rep.north);
\draw[flow] (fid.east) -- ([xshift=4mm]rep.south west);
\draw[flow] (harm.north) |- ($(ext.north)+(0,6mm)$) -| ([xshift=6mm]rep.north);
% declared edges whose upstream output the module does not read
\draw[flow, draw=hsGreyLine] (harm.east) -- node[flowlabel, above, sloped, pos=0.5]{order only} (ext.west);
\draw[flow, draw=hsGreyLine] (legit.east) -- node[flowlabel, below, sloped, pos=0.5]{order only} (fid.west);
% legend
\node[desc, text width=34mm, align=left] at (88mm,-36mm) {Teal roots read facts only.\\Blue modules read one root.\\Repair averages three.};
\end{tikzpicture}

\caption{The default EM-DAG as configured in \code{default.yaml}. Each panel names a module and
the moral-vector dimension it writes, and each circle gives its place in the evaluation order. Two
declared edges, drawn grey, fix the order without passing any data.}
\label{fig:compiler-emdag}
\end{figure}
% src: src/erisml_compiler/em_dag/profiles/default.yaml:11-24
% src: src/erisml_compiler/em_dag/modules/*.py (dependencies tuples)
% src: docs/architecture.md:236-249 (same edges)

The order that Kahn's algorithm produces for this profile is harm, rights, fairness, legitimacy,
epistemic, externality, care, autonomy, fidelity, repair. The emitted hardware sources record the
same order in their header.
% src: src/erisml_compiler/silicon/examples/erisml_em_dag.cpp:5

\subsection{What each module computes}

Most modules share two aggregation helpers. Let $F$ be the set of facts a module selects, and let
each fact $f$ have a severity weight $w(f)$ and a confidence $c_f$. The severity weights are
$0.2$ for minor, $0.5$ for moderate, $0.85$ for grave and $1.0$ for catastrophic, with $0.5$ when
the severity is missing. The negative aggregate is
\begin{equation}
v = -\max_{f \in F} w(f), \qquad c = \frac{1}{|F|}\sum_{f \in F} c_f, \qquad u = 1 - c,
\label{eq:compiler-aggneg}
\end{equation}
and the positive aggregate has the same form with the sign of $v$ reversed. With $F$ empty, both
return $v = 0$, $c = 1$, $u = 0$ and a neutral direction. The aggregate takes the worst fact, not
a sum, so several moderate facts never add up to a grave one.
% src: src/erisml_compiler/em_dag/modules/_helpers.py:35-40 (SEVERITY_WEIGHT)
% src: src/erisml_compiler/em_dag/modules/_helpers.py:93-165 (severity_score, aggregate_negative, aggregate_positive)

The five roots select facts by kind and, for most, by keywords in the fact's description.
\begin{itemize}
\item \emph{Harm} takes every fact of kind \code{harm} or \code{non\_maleficence} and applies the
negative aggregate.
\item \emph{Rights} takes harm facts whose description contains the string \code{right}.
\item \emph{Fairness} takes facts of kind \code{justice} whose description mentions unfairness,
bias or discrimination.
\item \emph{Legitimacy} takes legitimacy facts whose description mentions void, tyrannical,
coercive or fraudulent authority, together with every coercion fact.
\item \emph{Epistemic} takes truth and deception facts whose description mentions deception,
lying, withholding or manipulation. It records the deception and leaves to the verdict layer
whether the deception was justified.
\end{itemize}
% src: src/erisml_compiler/em_dag/modules/harm.py:10-24
% src: src/erisml_compiler/em_dag/modules/rights.py:10-28
% src: src/erisml_compiler/em_dag/modules/fairness.py:10-32
% src: src/erisml_compiler/em_dag/modules/legitimacy.py:10-36
% src: src/erisml_compiler/em_dag/modules/epistemic.py:10-38

The second tier reads upstream outputs or upstream concepts. The \emph{autonomy} module applies the
negative aggregate to consent and coercion facts whose description signals coercion or imposition.
It then composes with legitimacy. Let $s_A$ be its own score and $s_L$ the legitimacy score. If
$v_L < -0.5$ and $v_A > v_L$, the module replaces its score with
\begin{equation}
v_A' = v_L, \qquad c_A' = \min(c_A, c_L), \qquad u_A' = \max(u_A, u_L),
\label{eq:compiler-autonomy}
\end{equation}
with a negative direction. The code comment gives the reason. Consent under illegitimate authority
is void, so autonomy cannot score better than a severely compromised legitimacy.
% src: src/erisml_compiler/em_dag/modules/autonomy.py:1-55

The \emph{fidelity} module reads commitments. Let $V$ be the violated commitments and $A$ the
active, defeasible or fulfilled ones. Its value is
\begin{equation}
v_{\mathrm{fid}} =
\begin{cases}
-0.8 & \text{if } V \neq \emptyset,\\
0.6 & \text{else if some commitment in } A \text{ is active but defeasible},\\
0.85 & \text{else if } A \neq \emptyset,\\
0 & \text{otherwise,}
\end{cases}
\label{eq:compiler-fidelity}
\end{equation}
with confidences $0.9$, $0.85$, $0.85$ and $1$. The module declares a dependency on legitimacy,
and the default code does not read the legitimacy output. The dependency fixes the evaluation
order.
% src: src/erisml_compiler/em_dag/modules/fidelity.py:18-86

The \emph{externality} module takes externality facts whose subjects include a non-consenting
third party. If none qualify, it takes all externality facts. The set of non-consenting third
parties comes from a helper described below. The module declares a dependency on harm and does not
read the harm output either.
% src: src/erisml_compiler/em_dag/modules/externality.py:1-44

The \emph{care} module is the positive counterpart of harm. With care facts present it applies
the positive aggregate. Without them, if the graph holds at least one stakeholder of high or
extreme vulnerability and at least one stakeholder with the protector role, it returns
$v = 0.7$ with confidence $0.75$. Otherwise it returns zero.
% src: src/erisml_compiler/em_dag/modules/care.py:20-61

The \emph{repair} module measures the moral debt left after an act. Let
$N = \{\,d \in \{\mathrm{harm}, \mathrm{externality}, \mathrm{fidelity}\} : v_d < 0\,\}$ be the
upstream modules with negative scores. Then
\begin{equation}
v_{\mathrm{rep}} =
\begin{cases}
-\dfrac{1}{|N|}\displaystyle\sum_{d \in N} |v_d| & \text{if } N \neq \emptyset,\\[2ex]
0 & \text{otherwise,}
\end{cases}
\label{eq:compiler-repair}
\end{equation}
with confidence $0.8$ in the first case. Its explanation names the violated commitments. Repair is
the only module that averages, and it averages only the negative upstream scores.
% src: src/erisml_compiler/em_dag/modules/repair.py:17-73

\subsection{The graph-native helpers}

The modules do not read the representation's lists directly. They call six helpers.
\begin{itemize}
\item \code{facts\_of\_kind} returns the facts of one kind.
\item \code{active\_commitments} and \code{violated\_commitments} return commitments by status.
\item \code{vulnerable\_stakeholders} and \code{stakeholders\_with\_role} return stakeholders.
\item \code{nonconsenting\_third\_party\_ids} returns the identifiers of parties who bear an act
without consent.
\end{itemize}
The architecture document names five of them. Each helper queries the typed graph when one is attached and reads
the flat lists otherwise. Fact nodes carry their kind as a label, so \code{facts\_of\_kind} filters
fact nodes by label and rebuilds each fact from its payload.
% src: src/erisml_compiler/em_dag/modules/_helpers.py:8-21, 43-90
% src: src/erisml_compiler/ir/graph/promote.py:158-161 (labels=[kind])
% src: docs/architecture.md:256-265

The helper for non-consenting third parties shows what the graph adds. On the graph path it
returns the stakeholders labelled \code{nonconsenting\_third\_party}, together with
\begin{equation}
T = \{\, t : (a \xrightarrow{\code{imposes\_on}} t) \in E \ \wedge\ (t \xrightarrow{\code{consents\_to}} a) \notin E \,\},
\label{eq:compiler-thirdparty}
\end{equation}
the targets of an imposition edge with no paired consent edge. This is a typed-edge reading of
bearing an act without consent. The flat path instead matches role labels and a consent status of
not obtained or coerced. The project holds a golden test that compares the module outputs on the
graph path against a baseline captured on the flat path, for each of the three bundled scenarios.
% src: src/erisml_compiler/em_dag/modules/_helpers.py:210-243
% src: tests/test_em_dag_graph_native.py:9, 34, 57-78 (golden_em_dag_flat.json)

\subsection{From module outputs to the moral vector, and the ethos profile}

The moral vector has ten dimensions, one per module in the default profile. The projection copies
each module's score into the dimension that module declares. A dimension that no module owns gets
a neutral zero with an explanation saying so.
% src: src/erisml_compiler/ir/schemas.py:250-288 (MORAL_DIMENSIONS, MoralVector)
% src: src/erisml_compiler/evaluation/moral_vector.py:32-69

An ethos profile scales module values at this step. For module $m$ with weight $\omega_m$, the
projected value is
\begin{equation}
v_m' = \min\bigl(1, \max(-1, \omega_m v_m)\bigr),
\label{eq:compiler-ethos}
\end{equation}
with $\omega_m = 1$ for any module the profile does not list. Confidence and direction are left
unchanged, because the weight expresses salience and not the module's certainty.
% src: src/erisml_compiler/evaluation/moral_vector.py:19-55

Two fitted profiles ship, fit from the Dear Abby and AITA slices of Social Chemistry
101~\cite{forbes2020socialchem}. The Dear Abby profile lists weights for six modules. Harm and care
receive $1.74772$, legitimacy $0.970297$, fidelity $0.723662$, fairness $0.59368$ and rights
$0.216921$. Autonomy, epistemic, externality and repair are absent and so receive the default
weight of one. Each profile also carries per-module priors and coverage, a corpus fingerprint with
a SHA-256, a license and a citation, the hand-written mapping from moral foundations to modules,
and a list of bias notes. The notes state, for example, that the five moral foundations give no
signal to four of the ten modules. The loader keeps all of this text, and its docstring asks
callers to show it wherever the profile is used.
% src: docs/architecture.md:287-310
% src: src/erisml_compiler/em_dag/profiles/dear_abby_socialchem_v0.1.yaml:1-50, 76-103
% src: src/erisml_compiler/em_dag/dag.py:187-211 (load_ethos_profile)

\section{The finite-state layer and the silicon target}\label{sec:compiler-fsm}

\subsection{Three deterministic machines track moral state}

Tier~1 carries three finite-state machines. Each is a table from state and event tag to next
state. An event tag with no entry leaves the state unchanged. Terminal states absorb every event.
Each machine keeps a history of time index, state and event.
% src: src/erisml_compiler/fsm/__init__.py:1-16
% src: src/erisml_compiler/fsm/commitment_fsm.py:96-116

The commitment machine has seven states. From \code{active} a commitment can become
\code{active\_but\_defeasible}, \code{fulfilled}, \code{violated}, \code{void} when legitimacy
collapses, or \code{expired}. From the defeasible state it can return to \code{active}, become
\code{defeated}, or reach any of the other terminal states. The five terminal states are
\code{defeated}, \code{fulfilled}, \code{violated}, \code{void} and \code{expired}. The machine's
docstring lists the defeasibility conditions as a coerced vow, a vow to commit a wrong, a
catastrophic non-consensual externality and a conflict with a higher duty.
% src: src/erisml_compiler/fsm/commitment_fsm.py:1-93

The legitimacy machine has six states. A fully legitimate authority becomes defeasible on a
procedural violation, coercive on detected coercion and fraudulent when evidence is revealed. A
defeasible authority can be restored. A coercive authority escalates to tyrannical, and coercive
or tyrannical authority becomes void on catastrophic intent. A fraudulent authority becomes void
when evidence is revealed. Only \code{void} is terminal. The consent machine has four states. Consent
not obtained becomes \code{obtained} or \code{coerced}. Obtained consent can be withdrawn or
coerced on revisit. The states \code{coerced} and \code{withdrawn} are terminal.
% src: src/erisml_compiler/fsm/legitimacy_fsm.py:1-75
% src: src/erisml_compiler/fsm/consent_fsm.py:1-38

Every machine has at most seven states, so each fits a three-bit register. The architecture
document bounds the whole moral state by
$O(n_{\mathrm{commitments}} + n_{\mathrm{authorities}} + n_{\mathrm{consenters}})$ registers.
Figure~\ref{fig:compiler-commitment-fsm} draws the commitment machine.

\begin{figure}[t]
\centering
% Chapter 5. The commitment lifecycle machine (Tier 1), with the five absorbing states.
% src: C:\source\erisml-compiler-monitor\src\erisml_compiler\fsm\commitment_fsm.py:55-93
% src: C:\source\erisml-compiler-monitor\src\erisml_compiler\silicon\hls_emit.py:37-50 (ap_uint<3>, ap_uint<4>)
\begin{tikzpicture}[house]
\node[stratum, minimum width=26mm] (act) at (0,0) {active};
\node[stratum, minimum width=26mm] (def) at (0,-36mm) {active but defeasible};
\node[absorbing, minimum width=22mm] (ful) at (70mm,12mm) {fulfilled};
\node[absorbing, minimum width=22mm] (vio) at (70mm,0mm) {violated};
\node[absorbing, minimum width=22mm] (voi) at (70mm,-12mm) {void};
\node[absorbing, minimum width=22mm] (exp) at (70mm,-24mm) {expired};
\node[absorbing, minimum width=22mm] (dft) at (70mm,-38mm) {defeated};
% between the two live states
\draw[flow] ([xshift=-5mm]act.south) -- node[flowlabel, left, align=right]{defeasibility\\condition\\triggered} ([xshift=-5mm]def.north);
\draw[flow] ([xshift=5mm]def.north) -- node[flowlabel, right, align=left]{resolved\\in favour} ([xshift=5mm]act.south);
% from active
\draw[flow] (act.east) -- node[flowlabel, above, pos=0.8, sloped]{fulfilling} (ful.west);
\draw[flow] (act.east) -- node[flowlabel, above, pos=0.8, sloped]{violating} (vio.west);
\draw[flow] (act.east) -- node[flowlabel, above, pos=0.8, sloped]{legitimacy collapses} (voi.west);
\draw[flow] (act.east) -- node[flowlabel, above, pos=0.8, sloped]{expiration} (exp.west);
% from defeasible, the same four tags plus the override
\draw[flow, draw=hsFlow!60] (def.east) -- (ful.west);
\draw[flow, draw=hsFlow!60] (def.east) -- (vio.west);
\draw[flow, draw=hsFlow!60] (def.east) -- (voi.west);
\draw[flow, draw=hsFlow!60] (def.east) -- (exp.west);
\draw[flow] (def.east) -- node[flowlabel, below, sloped]{defeasibility overrides} (dft.west);
% register note
\node[panel=Grey, text width=40mm, anchor=north west] at (-14mm,-48mm)
  {{\scriptsize state in an \texttt{ap\_uint<3>} register,\\event tag in \texttt{ap\_uint<4>}}};
\node[panel=Red, text width=40mm, anchor=north west] at (52mm,-48mm)
  {{\scriptsize terminal states absorb every event,\\an unknown tag keeps the state}};
\end{tikzpicture}

\caption{The commitment lifecycle machine. Two live states lead to five absorbing states by named
event tags, and the defeasible state reaches the same four outcomes as the active state plus
defeat.}
\label{fig:compiler-commitment-fsm}
\end{figure}
% src: docs/architecture.md:271-285
% src: src/erisml_compiler/fsm/commitment_fsm.py:44-46

\subsection{What \code{silicon-emit} produces}

The command \code{eris-compile silicon-emit --out-dir} writes four files. The first,
\code{erisml\_fsm.cpp}, holds the three machines. Each becomes a state enumeration of type
\code{ap\_uint<3>}, an event enumeration of type \code{ap\_uint<4>}, and a step function. The step
function returns early from any terminal state and otherwise switches on state and event. It
carries the Vitis HLS pragmas \code{INLINE} and \code{PIPELINE II=1}. The transition cases are
generated from the same Python tables the machines use, so the cases cannot drift from the Python
reference. The lists of state and event names that number the enumerations are written out in the
emitter by hand.
% src: src/erisml_compiler/cli.py:613-660 (cmd_silicon_emit)
% src: src/erisml_compiler/silicon/hls_emit.py:1-12, 21-23 (imports _TRANSITIONS)
% src: src/erisml_compiler/silicon/hls_emit.py:37-201 (enum, switch, emit_fsm_cpp)

The second file, \code{erisml\_em\_dag.cpp}, holds the EM-DAG. It loads the chosen profile,
declares a fixed-point score type \code{ap\_fixed<16,4>} by default, and emits one function per
module whose arguments are the feature structure and the outputs of the module's dependencies. A
top function marked \code{DATAFLOW} calls the modules in topological order. Each module body is a
placeholder that returns a score of zero and a confidence of 200 on a scale of 255. The file
header says so, and the command's help text calls the file a skeleton. The third file,
\code{erisml\_top.cpp}, is the kernel entry point with AXI interface pragmas. It calls the EM-DAG
pipeline and sets the packed verdict word to zero. A comment reserves its bit layout, three bits
for the verdict kind, eight for confidence and one for escalation. The fourth file is a Makefile
that calls \code{v++} for the Alveo U55C platform, with an \code{hw\_emu} target for hardware
emulation.
% src: src/erisml_compiler/silicon/hls_emit.py:209-320 (emit_em_dag_pipeline, placeholder bodies)
% src: src/erisml_compiler/silicon/hls_emit.py:328-376 (emit_top_module, packed_verdict = 0)
% src: src/erisml_compiler/silicon/hls_emit.py:379-419 (emit_makefile, PLATFORM xilinx_u55c_gen3x16_xdma_3_202210_1)

Fixed-point conversion uses a signed Q format with $W$ total bits and $I$ integer bits, so
$f = W - I$ fractional bits. A value $x$ quantizes to the integer
\begin{equation}
q(x) = \operatorname{clip}\bigl(\operatorname{round}(x \cdot 2^{f}),\ -2^{W-1},\ 2^{W-1}-1\bigr),
\label{eq:compiler-quant}
\end{equation}
which saturates rather than wraps. The defaults are $W = 16$ and $I = 4$, and the command accepts
other widths.
% src: src/erisml_compiler/silicon/fixed_point.py:65-125
% src: src/erisml_compiler/cli.py:627-628

\subsection{What the silicon path has and has not shown}

Figure~\ref{fig:compiler-silicon} gathers the emitter's inputs, its four output files and the
recorded status of each. The silicon-target document lists the status of each step. The strict subset for the machines and
the EM-DAG, the fixed-point port of the arithmetic helpers and the emitter are marked done. The
document states that continuous integration verifies the emitted sources build. The project's
tests check that the emitted file contains the three machines, has balanced braces and carries the
pragmas, and that the EM-DAG file is in topological order with a \code{DATAFLOW} pragma.
Bit-exact equivalence between the emitted code and the Python reference is marked not done. FPGA
synthesis and on-board validation are marked blocked on the bitstream pipeline. A formal
equivalence proof is future work.
% src: docs/silicon_target.md:107-117 (status table)
% src: tests/test_silicon.py:20-97

\begin{figure}[t]
\centering
% Chapter 5. What eris-compile silicon-emit produces and where each file stands.
% src: C:\source\erisml-compiler-monitor\src\erisml_compiler\cli.py:613-660
% src: C:\source\erisml-compiler-monitor\src\erisml_compiler\silicon\hls_emit.py:115-419
% src: C:\source\erisml-compiler-monitor\docs\silicon_target.md:107-117
\begin{tikzpicture}[house,
  file/.style={panel=#1, text width=38mm, minimum height=12mm}]
% sources
\node[panel=Teal, text width=30mm, minimum height=16mm] (fsmsrc) at (0,6mm)
  {\textcolor{hsTeal}{\textbf{FSM tables}}\\{\scriptsize three Python\\transition dicts}};
\node[panel=Teal, text width=30mm, minimum height=16mm] (dagsrc) at (0,-24mm)
  {\textcolor{hsTeal}{\textbf{EM-DAG profile}}\\{\scriptsize \texttt{default.yaml}\\topological order}};
% emitter
\node[panel=Purple, text width=26mm, minimum height=56mm] (emit) at (40mm,-10.5mm)
  {\textcolor{hsPurple}{\textbf{silicon-emit}}\\[2pt]{\scriptsize fixed point\\\texttt{ap\_fixed<16,4>}\\by default}};
\node[step=Purple] at (emit.north west) {1};
% emitted files
\node[file=Green] (f1) at (84mm,12mm) {\textcolor{hsGreen}{\textbf{erisml\_fsm.cpp}}\\{\scriptsize cases generated from the tables}};
\node[file=Amber] (f2) at (84mm,-3mm) {\textcolor{hsAmber}{\textbf{erisml\_em\_dag.cpp}}\\{\scriptsize skeleton, every score is 0}};
\node[file=Amber] (f3) at (84mm,-18mm) {\textcolor{hsAmber}{\textbf{erisml\_top.cpp}}\\{\scriptsize packed verdict set to 0}};
\node[file=Grey] (f4) at (84mm,-33mm) {\textcolor{hsGrey}{\textbf{Makefile}}\\{\scriptsize \texttt{v++} for the Alveo U55C}};
\node[step=Green] at (f1.north west) {2};
\draw[flow] (fsmsrc.east) -- (emit.west |- fsmsrc.east);
\draw[flow] (dagsrc.east) -- (emit.west |- dagsrc.east);
\foreach \f in {f1,f2,f3,f4} { \draw[flow] (emit.east |- \f.west) -- (\f.west); }
% status row
\node[panel=Grey, text width=118mm, anchor=north west] (status) at (-19mm,-48mm)
  {\textcolor{hsGrey}{\textbf{Status recorded in the silicon-target document}}\\[2pt]
   {\scriptsize emitted sources build \quad bit-exact equivalence with Python not done \quad FPGA synthesis blocked \quad formal equivalence future work}};
\node[step=Grey] at (status.north west) {3};
\end{tikzpicture}

\caption{The silicon emitter turns the Python machine tables and the EM-DAG profile into four
Vitis HLS files. Only the machine file carries working logic, and the EM-DAG and top-level files
are skeletons that return zero.}
\label{fig:compiler-silicon}
\end{figure}

Two facts follow from the code. The emitted EM-DAG has the topology of the profile but not its
arithmetic, so no score computed in hardware can yet match a score computed in Python. The
document's resource figures for the default ten-module graph, fewer than five thousand lookup
tables and a critical path of at most ten nanoseconds, are stated as expectations. They are not
measurements.
% src: src/erisml_compiler/silicon/hls_emit.py:221-224, 289-290
% src: docs/silicon_target.md:75-79 ("Expected resource usage")

\RESULT{Bit-exact equivalence of emitted C++ against the Python reference on the pipeline test
vectors, and synthesis resource and timing figures, once measured.}

The tier description names a Mahalanobis evaluator as part of Tier~1, and the fixed-point helpers
mention porting it. The source tree has no Mahalanobis implementation. The term appears only in
docstrings.
% src: src/erisml_compiler/tiers.py:7, 79
% src: src/erisml_compiler/silicon/fixed_point.py:3 ; grep -rni mahalanobis src/ returns docstrings only

The document also states what the silicon target is not. It does not replace the
natural-language front end, which stays in software. It gates proposed actions and does not
generate them. It does not answer whether a deployment is lawful.
% src: docs/silicon_target.md:38-47, 95-105

\section{The audit chain and its exports}\label{sec:compiler-audit}

\subsection{Two hashes anchor a compile}

The compiler computes two content hashes. The graph hash is the SHA-256 of a canonical JSON
encoding of the moral graph,
\begin{equation}
h_G = \operatorname{SHA256}\bigl(\operatorname{canon}(G)\bigr).
\label{eq:compiler-graphhash}
\end{equation}
The encoding sorts nodes by identifier, sorts each node's labels, sorts every payload's keys
recursively, and sorts edges by source, destination, kind and payload. Two graphs with the same
nodes and edges therefore hash the same regardless of insertion order. Edges that differ only in
payload count as distinct.
% src: src/erisml_compiler/ir/graph/canonical.py:1-82
% src: docs/architecture.md:106-112

The representation hash $h_{\mathrm{IR}}$ is the SHA-256 of the representation dumped to JSON
with sorted keys. It excludes the audit record itself, the document timestamp and the source file
path. The document identity stays covered through the text hash stored in the document record.
% src: src/erisml_compiler/audit/hash_chain.py:22-41

The graph promotion has a round-trip property. Converting a graph to flat lists and back gives
the same graph hash for every bundled example.
% src: docs/architecture.md:123-127

\subsection{The audit record}

Pass~12 builds the \code{AuditRecord}. It holds $h_{\mathrm{IR}}$, the hash of the source text, the
schema version \code{erisml\_compiler\_ir\_v0.3}, the compiler version, the tier, the extractor
name, an optional model version, the EM-DAG profile name, $h_G$, the ethos profile name with the
SHA-256 of its YAML file, a UTC timestamp and the list of pass records.
% src: src/erisml_compiler/audit/hash_chain.py:44-71
% src: src/erisml_compiler/ir/schemas.py:365-386
% src: src/erisml_compiler/__init__.py:12

The architecture document says that two compiles of the same case yield bit-identical audit
records. The code supports a narrower statement. The same input, extractor and profiles produce
the same $h_G$ and the same $h_{\mathrm{IR}}$. The audit record as a whole also holds a timestamp
and measured pass durations, which differ between runs.
% src: docs/architecture.md:224-226
% src: src/erisml_compiler/audit/hash_chain.py:69 (timestamp_utc)
% src: src/erisml_compiler/audit/provenance.py:15-25 (duration_ms)

The decision proof extends the chain after Pass~12. It is a dictionary that hashes the moral
tensor's values, shape, axes and three moral-bearing metadata keys, hashes the ethical facts and
hashes a profile record. It lists the forbidden and ranked options. Its \code{previous\_proof\_hash}
field defaults to the audit record's $h_{\mathrm{IR}}$, and its \code{proof\_hash} is the SHA-256
of every other field. Because the proof is built after the audit record, $h_{\mathrm{IR}}$ does not
cover it. The link runs from the proof back to the representation.
Figure~\ref{fig:compiler-audit} draws the chain.
% src: src/erisml_compiler/erisml_backend/v3_phase6.py:195-337
% src: src/erisml_compiler/pipeline/orchestrator.py:468-507 (audit before proof)

\begin{figure}[t]
\centering
% Chapter 5. The audit chain of one compile.
% src: C:\source\erisml-compiler-monitor\src\erisml_compiler\audit\hash_chain.py:18-71
% src: C:\source\erisml-compiler-monitor\src\erisml_compiler\ir\graph\canonical.py:48-82
% src: C:\source\erisml-compiler-monitor\src\erisml_compiler\erisml_backend\v3_phase6.py:195-337
% src: C:\source\erisml-compiler-monitor\src\erisml_compiler\pipeline\orchestrator.py:202-214 (ethos sha256)
\begin{tikzpicture}[house,
  src/.style={panel=#1, text width=30mm, minimum height=14mm}]
% inputs and their hashes
\node[src=Blue] (txt) at (0,0) {\textcolor{hsBlue}{\textbf{source text}}\\{\scriptsize SHA-256 of the raw text}};
\node[src=Teal] (gr) at (0,-19mm) {\textcolor{hsTeal}{\textbf{moral graph}}\\{\scriptsize $h_G$ over sorted nodes\\and edges}};
\node[src=Purple] (ir) at (0,-38mm) {\textcolor{hsPurple}{\textbf{representation}}\\{\scriptsize $h_{\mathrm{IR}}$ without audit,\\timestamp, file path}};
\node[src=Green] (eth) at (0,-57mm) {\textcolor{hsGreen}{\textbf{ethos profile}}\\{\scriptsize SHA-256 of the YAML file}};
% the audit record
\node[panel=Orange, text width=36mm, minimum height=66mm] (aud) at (48mm,-28.5mm) {};
\node[stepname=Orange, anchor=north] at ($(aud.north)+(0,-2mm)$) {Audit record};
\node[detail, text width=31mm, anchor=north] (d1) at ($(aud.north)+(0,-8mm)$) {text hash, $h_G$, $h_{\mathrm{IR}}$};
\node[detail, text width=31mm, below=2mm of d1] (d2) {ethos name and hash};
\node[detail, text width=31mm, below=2mm of d2] (d3) {versions, tier, extractor,\\EM-DAG profile};
\node[detail, text width=31mm, below=2mm of d3] (d4) {timestamp and pass\\durations, vary per run};
\node[step=Orange] at (aud.north west) {12};
% decision proof
\node[panel=Grey, text width=30mm, minimum height=30mm] (prf) at (92mm,-28.5mm)
  {\textcolor{hsGrey}{\textbf{Decision proof}}\\[2pt]{\scriptsize tensor, facts and\\profile hashes\\[1pt]previous\_proof\_hash\\$=h_{\mathrm{IR}}$\\[1pt]proof\_hash over\\every other field}};
\foreach \s in {txt,gr,ir,eth} { \draw[flow] (\s.east) -- (aud.west |- \s.east); }
\draw[flow] (aud.east) -- node[flowlabel, above]{built after} (prf.west);
\end{tikzpicture}

\caption{The audit chain of one compile. Four inputs are hashed into the audit record, and the
decision proof points back to the representation hash. The timestamp and the pass durations
change between any two compiles, and the hashes do not.}
\label{fig:compiler-audit}
\end{figure}

A separate function bundles an audit folder for a reviewer. The folder holds the representation
as JSON, the source text, the ErisML rendering and a readable report.
% src: src/erisml_compiler/audit/artifact.py:12-37

\subsection{The training export}

The RLEF export writes one JSON record per compile under schema \code{rlef\_v0.2}. It holds the
source text, the flat annotations, the moral-vector timeline, the DEME verdict, the module
outputs, the audit record and any human corrections. Since version 0.8.0 it also holds the moral
graph, with its hash, canonical JSON, nodes, edges and counts by kind, every projection result and
the disagreement record. A trainer can learn against the typed graph, against one theory's
verdict or against the timeline.
% src: src/erisml_compiler/export/rlef.py:1-81
% src: docs/architecture.md:335-372

\subsection{The I-EIP monitor interface}

The I-EIP monitor compares the text-side representation with probes on a model's internal
activations. It has three lenses. The text lens is the consequentialist projection's moral vector.
The activation lens is a set of calibrated per-layer probes on hidden states, read from a mock
source, a local Hugging Face model or a remote host over SSH. The delta lens compares the two and
runs an equivariance test under meaning-preserving rewrites.
% src: docs/i_eip_monitor.md:1-38
% src: docs/architecture.md:312-333
% src: src/erisml_compiler/monitor/__init__.py:1-46

The comparison scores each dimension $d$ with weight $\omega_d$. Let $\Delta_d$ be the activation
value minus the text value, and let $\delta_d = 1$ when the two directions differ. The divergence is
\begin{equation}
D = \frac{\sum_d \omega_d \min\bigl(1, |\Delta_d| + 0.5\,\delta_d\bigr)}{\sum_d \omega_d},
\label{eq:compiler-divergence}
\end{equation}
and the comparison flags for review when $D \geq 0.35$, when more than two directions differ, or
when a dimension uncertain on both sides moves by at least $0.25$. Six named failure modes set
\code{requires\_human\_review}. They are \code{text\_internal\_mismatch},
\code{layerwise\_drift}, \code{group\_symmetry\_break}, \code{probe\_uncertainty\_spike},
\code{audit\_chain\_break} and \code{rho\_non\_orthogonal}. The monitor's only authorised output on
failure is that flag with a report. It never overrules a projection's verdict. Chapter~12 reports
the monitor study.
% src: src/erisml_compiler/delta/compare.py:81-190 (compare_morals, defaults)
% src: src/erisml_compiler/delta/failure_modes.py:1-60
% src: docs/i_eip_monitor.md:5-8

The monitor runs outside the silicon path by design. Activation hooks need a runtime that can read
a transformer's hidden states. The monitor document places the silicon path in the real-time loop
and the monitor out of band on a sampled subset of inputs.
% src: docs/i_eip_monitor.md:118-125

\section{Trust boundaries and stated limits}\label{sec:compiler-boundaries}

\subsection{Six boundaries, each with a mediator}

The architecture names six trust boundaries. Table~\ref{tab:compiler-trust} lists them with the
mechanism that mediates each one.
% src: docs/architecture.md:415-424

\begin{table}[t]
\centering\footnotesize
\begin{tabular}{p{30mm}p{30mm}p{62mm}}
\hline
Boundary & Untrusted side & Mediator \\
\hline
Text input & user text & extractor and canonicalizer \\
LLM output (Tier 3) & language model & critic, canonicalizer and schema validators \\
Probe output (Tier 2.5) & probe checkpoint & calibration provenance with probe and corpus hashes, anchored to the monitor trace \\
Activation source & remote inference host & SSH with a literal harness string and no remote code update \\
Ethos profile & YAML on disk & YAML SHA-256 in the audit record \\
Audit chain & downstream consumer & $h_{\mathrm{IR}}$, $h_G$ and the decision-proof hash \\
\hline
\end{tabular}
\caption{The six trust boundaries of the compiler. In each row the compiler is the trusted side.
Two mediators filter what crosses, three record a hash of it, and one restricts what runs remotely.}
\label{tab:compiler-trust}
\end{table}
% src: docs/architecture.md:417-424

The mediators work in three ways. Two filter. Extraction with canonicalization, and the critic
with the schema validators, change or reject what crosses. Three record. The probe provenance,
the ethos hash and the audit hashes leave content unchanged and make a later substitution
detectable. One restricts. The activation harness is a literal string in the compiler's source, and
no remote code update runs at inference time. The monitor document places that trust boundary at
the SSH account. The monitor document adds its own threat
model, with probe poisoning, activation spoofing and a faulty rewrite catalogue as the three named
threats.
% src: docs/i_eip_monitor.md:40-113

\subsection{What the architecture says it is not}

The architecture document ends with four disclaimers. This chapter restates them because they
bound every claim above.

The compiler is not neutral across metaethical positions. The substrate's categories,
stakeholders, maxims, commitments and ethical facts, are themselves choices. The split into
substrate and projections made that commitment smaller than it was before version 0.8.0. It did
not remove it.

The compiler is not a moral authority. It emits findings from several theories and declines to
choose between them when they disagree.

The compiler does not contain a complete Kantian, virtue or care analyser. The shipped findings
come from rule lists, a table from action kinds to virtues and role labels for dependency. The
universalizability gate now runs a solver, but over a hand-written model of nine facts.

The compiler does not replace human judgment. The monitor's only output on failure is a request
for human review. The projections produce findings meant to support deliberation.
% src: docs/architecture.md:426-447
% src: src/erisml_compiler/delta/universalizability_smt.py:46-58

\subsection{What the code adds to that list}

Reading the code adds five points to the document's own list. The rights, fairness, legitimacy,
epistemic and autonomy modules select facts by keyword matches on free-text descriptions, so their
scores depend on the extractor's wording. Two second-tier modules declare dependencies that their
code does not read, so in the default profile those edges fix order and carry no data. The polarity
map leaves three DEME verdict strings at \code{neutral}, so those verdicts never enter the
disagreement test. The emitted EM-DAG hardware is a skeleton. The audit record carries a timestamp
and timings, so only its two hashes are reproducible bit for bit.
% src: src/erisml_compiler/em_dag/modules/{rights,fairness,legitimacy,epistemic,autonomy}.py
% src: src/erisml_compiler/em_dag/modules/fidelity.py:21, externality.py:24
% src: src/erisml_compiler/projections/base.py:67-91
% src: src/erisml_compiler/silicon/hls_emit.py:221-224
% src: src/erisml_compiler/audit/hash_chain.py:69

The parts of the compiler that the home-care twin relies on are narrower than this chapter. The
twin loads its scene through the Tier~1 loader, canonicalizes with the registry canonicalizer, and
runs the scene runtime that Chapter~7 describes. The twin's code under \code{twin/} imports
neither the projections package nor the silicon package.
% src: C:\source\gtc-prototype\twin\brain.py:301-319
% src: grep -rn "projections|silicon|em_dag" C:\source\gtc-prototype\twin --include=*.py | grep import (no hits, 2026-10-03)
